% =====================================================================
% Слайд 4 — Данные: формат партитуры
% =====================================================================
\begin{frame}[fragile]{Данные: формат партитуры}
  \small
  Партитура хранится в JSON-формате \texttt{score.json}. Одна запись
  $=$ одно скрытое состояние HMM/HSMM (может быть аккордом).

  \begin{columns}[T,onlytextwidth]
    \column{0.45\textwidth}
\begin{verbatim}
{
  "index": 0,
  "pitches": [60, 64, 67],
  "nominal_onset":    2.999,
  "nominal_duration": 1.499
}
\end{verbatim}
    \begin{itemize}\footnotesize
      \item \texttt{generated\_dataset/} --- 4 синтетические пары;
      \item \texttt{midi/library/rach\_solo/} --- Рахманинов,
            \textbf{3791 state};
      \item для RL --- импортер ASAP (Foscarin 2020): 222 пьесы,
            1067 исполнений с note-level разметкой.
    \end{itemize}

    \column{0.55\textwidth}
    \centering
    % ГРАФИК 1 — Pitch vs time для rach_solo (схематически, ч/б).
% Данные сгенерированы как репрезентативная выборка из 3791 state
% Рахманинова — диапазон pitch 28-90, плотный кластер 50-80.
\begin{tikzpicture}
  \begin{axis}[
    width=6cm, height=4.5cm,
    xlabel={Время, c},
    ylabel={MIDI pitch},
    xmin=0, xmax=420,
    ymin=28, ymax=92,
    tick style={black},
    axis line style={black},
    grid=both,
    grid style={dotted, gray!50},
    label style={font=\footnotesize},
    tick label style={font=\scriptsize},
    title style={font=\footnotesize\itshape},
    title={Рахманинов: 3791 state партитуры}
  ]
    % Scatter point cloud — pseudo-distributed pitches
    \addplot+[
      only marks, mark=*, mark size=0.25pt,
      color=black, opacity=0.55,
    ] table[row sep=\\] {
      x y \\
      5  72\\  6 70\\  8 68\\ 10 65\\ 12 60\\ 14 64\\ 16 67\\ 18 72\\
      22 75\\ 24 79\\ 27 76\\ 30 72\\ 33 67\\ 36 64\\ 40 60\\ 43 64\\
      46 67\\ 49 70\\ 53 73\\ 57 76\\ 60 79\\ 63 75\\ 67 71\\ 71 67\\
      75 63\\ 79 59\\ 84 55\\ 88 51\\ 92 47\\ 96 51\\100 55\\104 60\\
     109 64\\114 68\\118 72\\123 75\\128 78\\133 75\\138 72\\143 68\\
     148 65\\153 60\\158 58\\163 55\\168 52\\172 48\\177 44\\181 41\\
     185 38\\189 44\\194 50\\199 55\\204 60\\209 65\\214 70\\220 75\\
     225 78\\230 80\\235 76\\240 72\\245 68\\250 64\\256 60\\262 56\\
     268 52\\274 56\\280 60\\286 65\\292 70\\298 74\\304 78\\310 81\\
     316 77\\322 73\\328 69\\335 65\\341 60\\347 56\\353 52\\359 48\\
     365 44\\372 49\\378 55\\384 60\\390 66\\396 71\\402 76\\408 80\\
      19 55\\ 21 58\\ 35 50\\ 38 56\\ 52 70\\ 56 64\\ 70 58\\ 80 64\\
      90 70\\ 99 64\\105 70\\110 76\\116 64\\125 66\\135 58\\150 70\\
     160 64\\180 80\\190 72\\200 68\\210 76\\230 88\\240 80\\260 72\\
     280 84\\300 76\\320 64\\340 72\\360 66\\380 75\\400 70\\
      15 36\\ 25 40\\ 35 44\\ 50 36\\ 80 40\\110 36\\140 32\\170 36\\
     200 40\\240 36\\280 32\\320 36\\360 40\\410 36\\
      11 80\\ 18 84\\ 33 88\\ 60 82\\ 95 90\\150 86\\200 88\\280 90\\
     350 84\\
    };
  \end{axis}
\end{tikzpicture}

    \figcaption{Рахманинов, 1-я часть концерта №\,2. Каждая точка~---
             состояние партитуры (верхняя нота аккорда).}
  \end{columns}
\end{frame}

